\chapter{Cross-lane operations}\label{ch:cross-lane-operations}

Cross-lane instructions move a value between the 32 lanes of one simdgroup without memory. Shuffles and
broadcasts read another lane's register; the native reductions combine all 32 (or 4) lanes in one instruction.

\leadhead{Instructions.}

\begin{itemize}
\item \op{14169} and \op{14170} implement the 32-lane Metal \code{simd\_shuffle\_xor}, and \op{13944} and \op{13945}
  the 4-lane forms.
\item \op{14157} reads a constant lane; \op{14158} reads the
  lane named in a register.
\item \op{16842} and \op{16858} implement Metal \code{simd\_sum} and \code{simd\_max} on float. Apple's \code{quad\_sum}
  compiles to the 10-byte form of \op{13857} (no glossary entry); \op{13856} is reachable from source but absent from the
  corpus (MM~\mm{25.82}).
\item \op{14208} writes the mask of active lanes (measured, glossary). \code{SR\_PVSIMD} and \code{SR\_TVSIMD}, the vote registers, are
  in \cref{ch:special-registers}.
\end{itemize}

\section{Shuffle-xor mask range}\label{sec:shuffle-xor-bounded-masks}

For an immediate mask m below 32, \op{14169} gives every lane L (0--31) the source value of lane L \^{} m, 32 of 32
lanes for each mask. For m of 32 or more every lane receives 0.0; the mask is not taken modulo 32. \op{13944} is bounded the same way to four lanes and returns 0.0 for any mask of 4
or more (measured, MM~\mm{12}, MM~\mm{17} rule 25). The glossary still tags these forms "compiled", but MM~\mm{12} measured them.

With a runtime mask, \op{13945} follows each lane's own mask exactly and \op{14170} does not, so the 32-lane form's
mask must be lane-uniform (measured, MM~\mm{12}). The SIMD runtime form appears nowhere in Apple's 5,010,028-instruction corpus, although Apple's compiler emits its 10-byte form for a
variable \code{simd\_shuffle\_xor} (MM~\mm{25.73}, MM~\mm{25.82}).

Below Metal, a 52-byte program authored by this compiler loads A{[}lane{]}, applies the immediate-mask shuffle-xor with mask 1 and stores C{[}lane{]},
with a waiting copy between the load and the shuffle. Under two input profiles with byte-identical code (normal fp32 patterns 0x3f800000 + 256 $\times$ lane,
and the small words 11 $\times$ lane + 5, which are fp32 subnormals), all 64 words came back as the bit pattern of lane L \^{} 1.
The subnormal patterns crossed lanes unchanged, so the shuffle does not apply the fp32 ALU's flush (\cref{sec:fp32-arithmetic};
measured, \code{docs/g17-first-dispatch-trials.md}, "SIMD cross-lane shuffle below Metal").

\section{Late operands}\label{sec:cross-lane-operations-consumers}

\code{simd\_shuffle\_xor} and \code{simd\_broadcast} read a loaded value before it arrived when they consumed it directly (measured, MM~\mm{25.90}).
In a two- and four-simdgroup combine that should have been exact, each arm returned one wrong value, a max below every
simdgroup's own max; the first butterfly step's shuffle had read a load that had not landed. A waiting copy before the
first shuffle made all four arms exact (\cref{sec:scoreboard}), and this compiler inserts that copy.

\section{Broadcast and inactive lanes}\label{sec:broadcast-inactive-lanes}

\op{14157} reads lane 0's value; it is measured only with every lane active, so lane 0
and "the first active lane" are not separated there (measured, MM~\mm{25.145.1}). A predicated region separates them: with lane 0 inactive
(first active lane 1), every active lane received 0, and \code{fma(x, x, broadcast(x))} returned exactly x*x, the broadcast
contributing 0, on 11 of 11 lanes that differed from the reference (measured, MM~\mm{25.144.6}). The constant-lane broadcast
therefore implements \code{simd\_broadcast(x, 0)} only.

Apple's compiler lowers a divergent \code{simd\_broadcast\_first} as four instructions (decoded; the roles are inferred, not measured
one by one, MM~\mm{25.144.6}): a \op{14060} of register 38 (\code{SR\_PVSIMD}, \cref{ch:special-registers}), \op{14214} over it (likely selecting the first active lane), \op{9989} turning
that into a lane index, and \op{14158} from that lane. This compiler does not lower that
sequence: it refuses a broadcast inside a predicated region and keeps the constant-lane broadcast for uniform code. No shuffle or vote in a partially active simdgroup is
established, and the emulator refuses them as well (MM~\mm{25.145.1}).

In this compiler, the constant-lane broadcast's source lifetime is operand 3. A template that left it at release made
\code{y = broadcast(x); z = fma(x, x, y)} read a released x, and those lanes stored 0 on hardware (MM~\mm{25.144.6}). The
release checker (release rules in \cref{ch:register-file}) treats the value a shuffle fetches from its source lane (operand 2
of the register-lane shuffle) as another lane's register, not a read of this lane's (\code{agxforge/g17/releasecheck.py}).

\section{Native reductions}\label{sec:native-reductions}

\code{simd\_sum}, \code{quad\_sum} and their product forms are a balanced adjacent-lane binary tree, every operation rounded to the
element type, exchanging over lane bits 0 to 4 in order; in half precision nothing accumulates at higher precision
(measured, MM~\mm{12}). For fp32, \op{16842} was tested on 64 rounding-sensitive 32-lane vectors with values spread
over $\pm 2^{20}$: it matches the adjacent pairwise tree, which is the xor butterfly in mask order 1, 2, 4, 8, 16, on 64 of 64. It
differs from both the project's earlier butterfly order (1, 8, 2, 4, 16) and an exact-then-rounded sum (measured, MM~\mm{25.141.16}). \op{16858} is exact on 32 lanes against an order-free reference (MM~\mm{25.90}).

As a half-precision example, lane 0 holds 2048.0 and every other lane 1.0 (measured, MM~\mm{12}; the arithmetic
below reproduces the result). Half has spacing 2 in {[}2048, 4096). \Cref{tab:half-sum-tree} follows lane 0's partial
sum up the tree.

\begin{xltabular}{\linewidth}{@{}>{\hsize=0.59\hsize}X>{\hsize=1.20\hsize}Xl>{\hsize=1.21\hsize}X@{}}
\caption{Half \code{simd\_sum} of 2048 and thirty-one ones, by tree level.}\label{tab:half-sum-tree}\\
\toprule
Tree level (lane bit) & Lane-0 partial & Exact & Note \\
\midrule
\endfirsthead
\tablecontinued{4}
\toprule
Tree level (lane bit) & Lane-0 partial & Exact & Note \\
\midrule
\endhead
0 & 2048 + 1 = 2049 $\to$ \textbf{2048} & 2049 & a tie; RNE keeps the even significand \\*
1 & 2048 + 2 = 2050 & 2051 & other partials are now 2 \\*
2 & 2050 + 4 = 2054 & 2055 &  \\*
3 & 2054 + 8 = 2062 & 2063 &  \\*
4 & 2062 + 16 = \textbf{2078} & 2079 & \code{simd\_sum} returns 2078.0 \\*
\bottomrule
\end{xltabular}

\code{quad\_sum} stops after lane bit 1 and returns 2050.0. A correctly rounded sum would give 2080 and 2052 (2079 and 2051 are
ties that round to even). The 1 lost at the first level is not recovered, because every later partial is computed and
rounded in half.

Reductions across simdgroups go through threadgroup memory: each simdgroup writes its value to threadgroup word sg, one
threadgroup barrier follows (\cref{ch:synchronization}), and every lane folds words 0 to S-1 in ascending order, so the result does not depend on arrival
order. The combine is exact at S = 2 (sum and max, 64 of 64 lanes) and S = 4 (128 of 128); a control with no exchange matched 0 lanes for sum
(measured, MM~\mm{25.90}).

The library tile reductions (\code{reduce\_rows}, \code{reduce\_columns} on tensor fragments) use a per-lane sequential chain and then a
butterfly over the reduction axis's lane bits, with identities and NaN handling of their own (\cref{sec:reductions}).

\section{Cross-lane dfdx}\label{sec:other-cross-lane-forms}

\op{1062} is decoded as a cross-lane \code{dfdx} with saturation. Its records ran on one lane only, so that function is
unconfirmed (MM~\mm{25.90}).

\textbf{Not known.} The behavior of the SIMD runtime-mask
shuffle-xor with a non-uniform mask; the roles of the four instructions in Apple's divergent broadcast; the reduction tree for half \code{simd\_sum} products and
for int types; cross-lane latency.

\textbf{Evidence.} MM~\mm{12}, MM~\mm{17} rule 25, MM~\mm{25.73}, MM~\mm{25.82}, MM~\mm{25.90}, MM~\mm{25.141.16}, MM~\mm{25.144.6}, MM~\mm{25.145.1};
\code{docs/g17-first-dispatch-trials.md}, \code{agxforge/g17/releasecheck.py}.
